% COMPUTATIONAL METHODS.  Since 2026-08-02 this is the whole of Sec. 2: the former "Setup and
% theory" section was folded in here, because the composed predictor, the error split and the
% estimator that measures it are the method, and JCIM does not carry a separate theory section.
% WHAT MAY NOT LEAVE THIS SECTION, however tempting the page saving -- a reader of the main text
% alone must never be able to conclude that a control was not run:
%   * the data, the scaffold split, and what the labels are;
%   * the closure implementation AND its validation, including where the validation stops (the
%     dispersion term), because the fidelity result of Sec. 3.2.3 rests on that boundary;
%   * the evaluation protocol: the three oracle applications, which results are three-seed and which
%     single-seed, and on what row set;
%   * the estimator cell, the convention rule, the stratification rule and the admissibility rule,
%     all four fixed here, before any margin is reported;
%   * every disclosure that qualifies a result -- the separately tuned arms, the scaffold-leak guard
%     being a property of the builder and not of these runs, the wrong-split sigma-stream build, and
%     the placeholder repository.
% Construction detail (the solver's damped map, implicit-function gradients, the three phases by
% name, the loss taxonomy and per-phase weights, the architecture, the NRTL parametrisation, the
% COSMO-SAC constants and the 51-bin grid, all hyperparameters) is in App. A; split hazards, stream
% provenance and the full leak audit in App. F.
\section{Computational methods}
\label{sec:methods}

\subsection{The composed predictor and its error split}
\label{sec:framework}
For a solute~$s$ in solvent~$v$ at temperature~$T$, the target is $y=\lnx$, the
log saturated mole fraction. Dissolving a solid trades two costs against each
other---breaking the crystal lattice ($\Phicry$) and mixing the freed molecules
against non-ideal solute--solvent interactions ($\lng$)---and solid--liquid
equilibrium (SLE) sets the log-solubility to their negative sum:
\begin{equation}
  \lnx = -\underbrace{\frac{\dHfus}{\Rgas}\!\Big(\tfrac1T-\tfrac1\Tm\Big)}_{\Phicry(T)\ \text{(crystal)}}
         \;-\;\underbrace{\lng(x_2,T)}_{\text{activity}} ,
  \label{eq:sle}
\end{equation}
with the constant-$\Delta C_p$ (heat-capacity-of-fusion) correction omitted for clarity. The
activity term is produced by a \emph{differentiable closure} $g$: a graph encoder predicts a
$\sigma$-profile---the distribution of screening (surface) charge density over a molecule,
the histogram a COSMO model consumes---for each molecule, and COSMO-SAC \cite{lin2002}, a
segment-activity reformulation of COSMO-RS, maps the pair of profiles to $\lng$. Inside that
map the segment-interaction energy---the electrostatic-misfit and hydrogen-bonding terms
setting what two surface segments cost each other---is the closure's \emph{kernel}, the part
the published revisions rewrite, so a closure variant is named for the year of its kernel:
the deployed \emph{2002 kernel}, and the \emph{2010 kernel} that types segments as
hydrogen-bond donors or acceptors instead of treating hydrogen bonding with one
parameter. Write $\zstar$ for the \emph{reference} physical latent, the
$\sigma$-profile tabulated in the VT-2005 database \cite{mullins2006}, and $g(\zstar)$ for the
closure output when fed it; the model's own head instead predicts a \emph{learned} profile $\hat z$
and evaluates $g(\hat z)$. Feeding $\zstar$ in place of $\hat z$ (a $\sigma$-oracle) isolates the
closure from encoder error: if the closure were well specified, $g(\zstar)$ should be the best
available activity estimate. Each trained configuration, and each closure variant evaluated, is an
\emph{arm}; a \emph{DirectGNN} control arm drops the closure and emits $\lnx$ directly
(Fig.~\ref{fig:arch}), and the two were tuned separately, so the gap between them is not
attributable to the bottleneck alone (\S\ref{sec:eval}; per-arm settings in
\S\ref{sec:si-methods}).

Everything below is stated for a general \emph{composed predictor} $\hat y=g(h(x))$ in which a
fixed closure $g$ consumes a learned physical intermediate $z=h(x)$, against a matched free
predictor that shares the encoder and drops $g$; solubility is one instance, \S\ref{sec:pka} a
second (pKa through a fixed Hammett closure) and \S\ref{sec:dial} a controlled synthetic family.
Nothing in Lemma~\ref{prop:decomp} or in \S\ref{sec:si-proofs} is specific to COSMO-SAC.

The crystal/activity split of Eq.~\eqref{eq:sle} is not identifiable from solubility alone, which
is why an unconstrained head can silently absorb crystal error into the activity term---an
indeterminate component end-to-end training leaves free rather than projecting out, closed only by
a rank-completing external constraint (Lemma~\ref{prop:ident}, \S\ref{sec:si-proofs}). That is
enabling background: it explains why an under-determined factor needs external single-component
data and why compensation between the factors is expected, and says nothing by itself about
whether grounding helps.

\subsubsection{The closure--insufficiency decomposition}
\label{sec:decomp}
Let $m$ be the activity target---the experimental infinite-dilution activity coefficient
$\lngi$---and $g(\zstar)$ the closure evaluated on the true latent.

\begin{lemma}[Exact closure--insufficiency decomposition and its one-sided bounds]\label{prop:decomp}
For a \emph{fixed} (non-fitted) closure $g$, with $\zstar$ denoting the closure's \emph{entire}
argument (the profile pair together with the temperature wherever the closure reads one),
\begin{equation}
\begin{aligned}
  \underbrace{\Ex\big[(m-g(\zstar))^2\big]}_{\text{ref.-input MSE}}
  &= \underbrace{\Ex\big[\Var(m\mid \zstar)\big]}_{\Binsuf\ (\text{irreducible})} \\
  &\quad + \underbrace{\Ex\big[(\Ex[m\mid\zstar]-g(\zstar))^2\big]}_{\Bclos\ (\text{decoder bias})} .
\end{aligned}
  \label{eq:decomp}
\end{equation}
The cross term vanishes identically, which is why $\zstar$ must be the
closure's whole argument and not a part of it (proof, \S\ref{sec:si-proofs}). The split is exact with no fitting of the closure
$g$, so the total and $\Bclos$'s Jensen lower bound are
fit-free. $\Binsuf=\Ex[\Var(m\mid\zstar)]$, however, must be estimated with a conditional-variance
smoother whose choice affects the verdict (\S\ref{sec:estimator}).
\end{lemma}
\noindent $\Bclos$ is the systematic discrepancy of a fixed physical map in the sense of
Kennedy--O'Hagan \cite{kennedy2001,brynjarsdottir2014}, and its dominance over $\Binsuf$ is the
empirical content of the claim that the accuracy ceiling is set by the closure rather than by the
inputs. Where model-discrepancy calibration \emph{fits} a correction from held-out observations,
here $g$ is fixed and $\zstar$ externally observed, so the map's error is \emph{attributed} rather
than corrected---available only where an external latent oracle exists, not in gray-box hybrids
(\S\ref{sec:si-positioning}). \textbf{Exactness is contingent on $g$ being fixed}: for a
\emph{fitted} decoder the cross term does not vanish and a plug-in point split is invalid, so we
report \emph{one-sided bounds} (proved in \S\ref{sec:si-proofs}):
\begin{itemize}
  \item $\Bclos \ge (\Ex[m]-\Ex[g])^2$ by Jensen: a constant offset between the
  decoder's mean output and the target is misspecification, not input insufficiency, and
  this bound needs \emph{no smoothness assumption}. Its one requirement is that $m$ and $g$
  share the $\lngi$ reference state, which the closure validation of \S\ref{sec:model} establishes.
  \item $\Binsuf \le \Ex[\Var(m\mid \mathrm{bin}(g(\zstar)))]$, a valid
  \emph{upper} bound (law of total variance, since a binning of $g(\zstar)$
  coarsens $\zstar$).
\end{itemize}
Because $\Binsuf$ is bounded above and $\Bclos$ below, a positive \emph{separation margin}
$\mathrm{MSE}-2\Binsuf^{\mathrm{up}}$ establishes that the closure's own error exceeds what input
insufficiency can account for, while a non-positive one is a failure to separate and licenses
nothing: a reversal would need a \emph{lower} bound on $\Binsuf$, which this instrument does not
supply.

\paragraph{What population, what is conditioned on, and what sets the resolution.}
$\Binsuf$ is well posed only once the population and the conditioning variable are named. The
population is the \emph{measurement superpopulation}: a draw is a nominal compound pair \emph{at a
temperature}, together with the conformer, tautomer and protonation state actually populated under
those conditions, and an independent experimental determination of $\lngi$ for it. The conditioning
variable is the closure's full argument, $(\zstar,T)$ where $g$ reads a temperature. Neither the
conformational state nor the replicate determination is part of it, so $\Var(m\mid\zstar,T)$ is the
spread of $\lngi$ across the conformational ensemble and across replicate determinations at one
reference profile and temperature, non-zero for physical reasons unrelated to sample size.
Experimental label noise is a \emph{component} of it, not a term beside it: an independent
zero-mean error on $m$ enters $\Var(m\mid\zstar,T)$ directly, so it is inside every bound we report
and never added on. What the design cannot do is \emph{observe} it, neither set holding two rows at
a common argument (\S\ref{sec:si-repro}), so $\Binsuf$ is reached only from above and $\Bclos$ only
from below. The bounds are conservative in one direction only---were the superpopulation variance
negligible the finding would be the stronger $\Bclos=\mathrm{MSE}$---so nothing reported can be an
artifact of over-estimating $\Binsuf$, and no entry is the share of the error the inputs
contribute. A low $P_{\mathrm{boot}}$, the cluster-bootstrap frequency that the measured margin
stays positive, is likewise loss of resolution and never evidence for the inputs; resolution is
governed by $\dim(\zstar)/n$.

\subsubsection{The decomposition and the grounding gap}
\label{sec:gap}
The decomposition is the measurement the conclusions rest on: assumption-light, needing neither a
representable encoder nor a lossless bottleneck. Alongside it we report a second, purely observable
quantity, the \emph{grounding gap} $\Ggap:=R_{\mathrm{orc}}-R_{\mathrm{e2e}}$, the oracle risk minus
the best end-to-end risk; the paradox is $\Ggap>0$. Under a representable encoder (A1) and a
lossless bottleneck (A2) it is sandwiched $0\le\Ggap\le\Bclos$, so $\Ggap>0$ would certify
$\Bclos>0$ (Theorem~\ref{thm:gap}, \S\ref{sec:si-proofs}). A1 is empirically doubtful---the encoder
is transfer-limited, a linear probe dropping from train $R^2\approx0.76$ to test
$R^2\approx0.45$---and A2 is unverifiable and physically dubious; when A2 fails, $\Ggap>0$ conflates
closure misspecification with input insufficiency, a well-specified closure with a lossy latent
giving $\Bclos=0$ yet $\Ggap>0$ (\S\ref{sec:dial}, and the pKa ortho sweep of \S\ref{sec:pka}). We
therefore read $\Ggap$ as corroborating evidence and base the model-dominance conclusion on the
decomposition.

\paragraph{The sign rule, and scope.}\label{sec:tax} Whether the \emph{trained} physics model
outperforms the \emph{trained} black box at sample size $n$ follows a sign rule with two halves: a
constrained prior helps only where the estimation variance it saves exceeds its asymptotic bias
$\Delta_\infty$, and a prior whose asymptotic bias is strictly positive always loses once the budget
passes a critical size, the variance it saves vanishing as data accumulate while $\Delta_\infty$
does not (Proposition~\ref{prop:tax}, \S\ref{sec:si-proofs}). Each physics-versus-control difference
below is a penalty at the budget actually trained, not an asymptotic one; there is no ``optimal
grounding'' theorem, and the decomposition does not improve out-of-distribution accuracy. Everything
here is measured at infinite dilution on two sets that differ in regime (\S\ref{sec:measure}), and
transfer of either to the finite-composition regime in which the paradox is observed is a conjecture
(\S\ref{sec:discussion}). We read these results against the model-discrepancy, sloppiness,
identifiability, misspecification, underspecification and null-space-prior literatures
(\S\ref{sec:si-positioning}). Spanning two chemical domains, this paper would otherwise use two
letters for more than one quantity each; Table~\ref{tab:notation} gives the renamings that separate
them, and every other symbol is defined where it is used.

% Notation table. Moved from the SI to the main text (2026-07-27) so that the symbols and the coined
% vocabulary are defined before the Results use them, and set at the end of the theory section
% (2026-07-29) when the separate notation-and-grading section was dissolved (criterion 7).
%
% SIMPLIFIED 2026-07-29 under criterion (8). It is a glossary, not a second methods section: one line
% per symbol, the qualification that belongs to a measurement kept where the measurement is made
% (Sec. 5.2, Tables 2-4, App. C). What left this table and where it now stands:
%   - the sets' sizes, temperature and heavy-atom ceiling, and the deposited file names -> Sec. 5.2(i)-(ii)
%   - "label noise is inside B_insuff, not beside it" -> Sec. 3.2 and Table 3, note k
%   - the margin's bootstrap-resolution rule -> the grading rule at the head of Sec. 5
%   - the full inventory of subscripted sigmas (Fisher design, cutoff, derivative) -> App. F, App. G
%   - "deflated" -> App. C, where the deflated estimators are tabulated
% Keep it to one line per row; a row that grows to three is a paragraph in the wrong place.
\begin{table*}[t]
\centering
\caption{Symbols and coined terms used throughout.}
\label{tab:notation}
{\footnotesize
\begin{tabular}{@{}lp{0.68\linewidth}@{}}
\toprule
Symbol / term & Meaning \\
\midrule
$g$ (closure), $h$ (encoder) & fixed, non-fitted physical map (COSMO-SAC; a Hammett linear free-energy relation, LFER, for pKa); learned graph network $x\mapsto z$ \\
$\zstar$, $\hat z$, $m$ & reference (oracle) vs.\ learned $\sigma$-profile; activity target (experimental $\lngi$). $\zstar$ is the closure's whole argument: the concatenated profile pair, written $(\zstar,T)$ where $g$ also reads a temperature \\
$\Bclos$, $\Binsuf$    & closure misspecification $\Ex[(\Ex[m\mid\zstar]-g(\zstar))^2]$; input insufficiency $\Ex[\Var(m\mid\zstar)]$, over the measurement superpopulation of \S\ref{sec:decomp} \\
separation margin      & $\mathrm{MSE}-2\Binsuf^{\mathrm{up}}$; positive witnesses $\Bclos>\Binsuf$ on the sample measured, negative loses that ordering rather than reversing it \\
the corner / the broad IDAC set & the two activity data sets the instrument is run on: the VT-2005-matchable pairs at one temperature, and the larger set matched to the University of Delaware (UD) $\sigma$-profile database across temperatures (\S\ref{sec:measure}) \\
LOTV; leakage-immune & law of total variance, the model-free binning bound on $\Binsuf$; \emph{leakage-immune}, fitting no model, so no held-out row borrows strength from a near-twin \\
\Ggap\ (grounding gap) & $R_{\mathrm{orc}}-R_{\mathrm{e2e}}$: oracle risk minus best end-to-end risk ($R_{\mathrm{free}}$ is the free black box's) \\
$\Gamma_{S}$; \hammett; $\sigma_\eta$ & COSMO-SAC segment activity coefficient (Eq.~\eqref{eq:cosmo}), \emph{not} the grounding gap \Ggap; Hammett substituent constant, renamed from $\sigma$; experimental label noise. Standing alone, $\sigma$ is the charge-density profile, and $\sigma(\cdot)$ with an argument is the generated $\sigma$-algebra \\
$\Delta_\infty$ (physics penalty) & asymptotic risk of the physics arm minus the budget-matched control \\
full / res             & COSMO-SAC with / without the Staverman--Guggenheim combinatorial term \\
compensating surrogate & a learned $\hat z$ that would \emph{cancel} the closure's error instead of matching $\zstar$; the hypothesis of \S\ref{sec:surrogate} \\
$P_{\mathrm{boot}}$    & clustered-bootstrap frequency that the \emph{measured} separation margin stays positive: the instrument's resolution on this design, not a posterior probability \\
\bottomrule
\end{tabular}}
\end{table*}


\subsection{Data, splits, and labels}
\label{sec:data}
Solubility measurements come from BigSolDB~2.0 \cite{krasnov2025bigsoldb}: $127{,}930$ rows over
$15{,}665$ solutes and $221$ solvents between $243$ and $438$\,K, partitioned by Bemis--Murcko
scaffold into $111{,}724$ training, $8{,}103$ validation and $8{,}103$ test rows so that every test
solute is unseen at training; a
melting-point-independent miscibility/structure filter sets row membership, and crystal $\Tm$,
$\dHfus$, Hansen parameters and infinite-dilution activity coefficients \lngi\ are merged on as
per-solute labels where available. Crystal supervision is scarce---the test and validation splits
contain essentially no measured $\dHfus$, the quantitative basis of the identifiability argument
(\S\ref{sec:framework}). The external single-component data used for grounding are an open
crystal-property pool ($\sim$15{,}000 melting points), the VT-2005 $\sigma$-profile database, and
ThermoML infinite-dilution activity coefficients \cite{frenkel2011thermoml}. Two hazards the
pipeline handles explicitly---melting-point unit detection, and the instability of the seeded
scaffold split across pipeline versions---are recorded in \S\ref{sec:si-repro}.

\paragraph{Disclosure: the scaffold guarantee is not certified for these runs.}
The released stream builder excludes any pool molecule whose scaffold appears in validation
or test and aborts on a missing split file --- but that is a property of the builder, not a
certified property of the runs reported here. One $\sigma$-stream build on disk was
generated against the wrong split files, $91$ of its $1425$ rows being held-out solutes by
canonical SMILES, and per-run stream provenance was not logged, so the artifacts cannot certify
which build fed which arm. \S\ref{sec:si-repro} gives the file-order evidence, the three facts that
bound what a leak could do, and the one molecule it touches in the $n{=}44$ analysis of
\S\ref{sec:surrogate}. We therefore do not assert the scaffold guarantee for these runs.

\subsection{Model, closure, and training}
\label{sec:model}
A \emph{shared} message-passing encoder embeds solute and solvent, a bipartite cross-attention block
lets them interact, and the resulting pair representation feeds every head; two interchangeable
encoder alternatives do not outperform it. A fusion head predicts $\Tm$ and $\dHfus$, giving the
ideal term $\Phicry(T)$ of Eq.~\eqref{eq:sle} with the constant-$\Delta C_p$ correction omitted
(audited in \S\ref{sec:si-mech}). The activity term \lng\ comes from one of three differentiable
closures. Two are controls: the fitted NRTL model and its symmetric one-parameter collapse, in which
the identifiability analysis is cleanest (\S\ref{sec:framework}). The third, parameter-free, is the
closure under study---a $\sigma$-profile head emitting a $51$-bin screening-charge-density profile
per molecule, and a COSMO-SAC layer mapping the profile pair to \lng\ with an optional
Staverman--Guggenheim combinatorial term, the full/residual-only distinction on which
Table~\ref{tab:claims} and \S\ref{sec:measure} turn. That closure over a network-\emph{predicted}
$\sigma$-profile, and the intervention that swaps it for the external reference, are the object of
study (\S\ref{sec:paradox}).

The saturated mole fraction solves $\lnx=-\Phicry(T)-\lng(x_2,T)$ as a damped fixed point in $x_2$
(\S\ref{sec:si-methods}); a bounded, gated ``adaptive physics correction'' perturbs the physical
parameters rather than the output, so the physical decomposition survives it. External
single-component data enter through \emph{grounding streams}---the training-time sense of grounding,
a separate operation from the evaluation-time substitution below. Each stream is a sidecar of
\emph{self-solvent} rows ($\mathrm{solvent}=\mathrm{solute}$, no solubility label, only the target
factor's mask active), interleaved into training under the guard above; the streams ground the
crystal factor from the melting-point pool and the activity factor from the $\sigma$-profile
database, identifying from external data a factor that solubility alone leaves non-identifiable
(\S\ref{sec:framework}) and that a black-box predictor cannot use. Training runs a three-stage
curriculum under a weighted loss: a warm-up on the auxiliary streams alone, then joint training
against solubility, then a low-rate stabilisation stage. The schedule pinned for the arms of
Fig.~\ref{fig:paradox} freezes the $\sigma$ head at the end of the warm-up, so the $\sigma$ stream
enters the warm-up only (\S\ref{sec:si-tables}).

\paragraph{The in-house closures, and where their validation stops.}
% Artifact: results/b_insuff/closure_reference_validation.json (2002, rmse 0.0035) via
% scripts/analysis/run_closure_reference_validation.py. The 2010 figure has no deposited
% artifact -- scripts/analysis/validate_cosmo2010_layer.py printed to stdout on the run it
% comes from -- which the SI states.
The dispersion-off closures in the fidelity measurement (\S\ref{sec:fidelity}) are our own
differentiable implementations, so those arms share one code path, and each consumes its own
model-prescribed $\sigma$-averaging (Mullins for 2002, Hsieh for 2010). Our 2002 layer---the
\emph{fair 2002} arm wherever the kernels are compared, because it reads all three profile
channels---reproduces the NIST reference COSMO-SAC-2002~\cite{bell2020cosmosac} to RMSE
$0.0035\,\ln\gamma$ on identical VT-2005 profiles, and our COSMO-SAC-2010/dsp
layer~\cite{hsieh2010,hsieh2014dsp}, built on the \emph{typed} latent of three
donor/acceptor-resolved $\sigma$-profiles per molecule in place of one, reproduces the NIST 2010
reference with dispersion \emph{off}. The dispersion term is \emph{not} validated to that standard,
which is why every dispersion-off contrast is computed on our layers while the \emph{2010,
dispersion on} runs are quoted from \texttt{cCOSMO}. Layer definitions, validation figures and the
full scope of the dispersion caveat are in \S\ref{sec:si-methods}.

\subsection{Evaluation protocol, and the DirectGNN control}
\label{sec:eval}
The \emph{oracle substitution} replaces the predicted $\sigma$-profile $\hat z$ with the reference
VT-2005 profile $\zstar$ before the closure, matched by canonical SMILES. It is applied three ways
to rule out implementation artifacts: at evaluation only; injected at training time so the model
co-adapts; and as a channel-swap control, injected at training under a coordinate-descent schedule
that freezes the crystal branch during the SLE phase, so that only the activity branch refits
against it. The last two are at seed $42$ only. An analogous override substitutes reference crystal
targets. The activity target for the decomposition (\S\ref{sec:measure}) is the infinite-dilution
$\lngi=\lim_{x_2\to0}\lng$, on two sets that differ in their reference database and their regime
(\S\ref{sec:broad}--\S\ref{sec:vt2005}).

DirectGNN shares the encoder, interaction, readout and pair representation, adds a thermometer
(cumulative-threshold) temperature encoding, and maps directly to \lnx\ with no heads, closure or
solver. The two arms therefore share their representation but not their inputs or their tuning: the
control receives an explicit temperature encoding the physics arm does not, seeing $T$ instead
through $\Phi(T)$ and the closure, and the optimisation settings were tuned per arm and differ. The
TGNNSolv$-$DirectGNN gap is accordingly a comparison of two separately tuned pipelines that differ
in the thermodynamic bottleneck, not an isolation of the bottleneck, and every physics-penalty
comparison in \S\ref{sec:results} carries that caveat.

The controlled comparisons of \S\ref{sec:paradox} are run at three seeds and reported as
mean$\pm$sd on a cross-arm row set; $5608$ is that lock, the test rows supervised and finite in
every arm, so all arms score identical rows. Single-seed, and labelled as such where each is
reported, are the training-time and channel-swap oracle controls, the output-residual arms
(\S\ref{sec:gateb}), the data-efficiency curve (\S\ref{sec:data-efficiency}) and the
loss-simplification ablation (\S\ref{sec:si-methods}). The two audit runs behind
\S\ref{sec:measure} are named there with the script and the artifact file of each; not every
artifact those scripts read is small enough to version, and the larger ones are deposited
separately (\S\ref{sec:discussion}).

\subsection{Estimating the split: one cell, one convention, one rule}
\label{sec:estimator}
Four analyst choices set the law-of-total-variance bound---the bin count, where the equal-count bin
edges fall, the within-bin variance convention, and the combinatorial convention---and all four are
fixed here, before any margin is reported.

\emph{The estimator cell.} Eight equal-count bins of $g(\zstar)$, the \emph{unbiased} within-bin
variance, and the row unit, applied identically to every set and every stratum whatever its size.
Enumerating every equal-count contiguous partition, the bin-edge choice is immaterial on the broad
IDAC set and material on the VT-2005-matched set, whose margin changes sign inside the envelope; we
accordingly read that set's $+0.05$ as consistent with zero and carry the claim on the broad set in
stratified form. The binning rule and its envelope are in \S\ref{sec:si-tables}.

\emph{The convention rule.} The margin depends on the combinatorial convention, so one rule governs
both sets: headline the \emph{deployed} residual-only convention, pairing $\mathrm{MSE}$ and the
bound computed under that same convention. Headlining the deployed convention on the smaller set and
the full convention on the broad one would be the favourable choice on each set taken separately;
one rule applied to both costs the broad set's margin about half its value. The Jensen
constant-offset bound is convention-specific---it separates under the full combinatorial convention
and inverts under the deployed residual-only default---so the separation reported in
\S\ref{sec:measure} rests on the model-free law-of-total-variance (LOTV) bound, which fits nothing
and so lets no held-out row borrow strength from a near-twin. $\Binsuf$ is one population quantity
while the two conventions coarsen two different statistics of it, so both bounds are valid on it,
and so is the smaller. The Jensen numbers are in \S\ref{sec:si-tables} and \S\ref{sec:si-proofs}.

\emph{The stratification rule.} The broad set is stratifiable where the VT-2005-matched set's $60$
pairs are not, on two axes that are pure functions of molecular structure: an ordered nine-class
cascade of substructure tests over the solvent, hydrogen-bonding role first and dominant functional
group second, and the solute's role as water or as an organic. The classifier reads a SMILES string
and nothing else---never $m$, $g$, the squared error, the source or the bootstrap---so no cell can
have been chosen on its outcome, which is a property of the code rather than an assurance. A stratum
is \emph{boundable} at $n\ge40$, where eight bins hold five rows; the smaller ones are reported at
that same cell rather than at a rule of their own, with the resample-adaptive rule
$\max(3,\lfloor n/10\rfloor)$ carried beside them in the deposited table, nothing being dropped for
being small.

\emph{The admissibility rule.} A stratum is \emph{admissible} only if it is \emph{(a)}~boundable at
that fixed cell and \emph{(b)}~robust to leave-one-source-out---its margin keeps its sign when each
source publication contributing to it is deleted in turn---in all four
unit\,$\times$\,convention cells, since the map reports all four. That last clause is load-bearing
and is never dropped: without it the set taken whole passes too. A stratum drawing on one publication
fails~(b) by non-testability: there is no deletion to perform, so the sign was never put at risk. A
deletion leaving fewer than $16$ rows leaves the fixed cell undefined and the sign unverifiable,
which fails~(b) as well. Everything else is \emph{reported in the map and not stated}, with its $n$,
its sources, its margin, its interval and the reason. The rule governs the cells that favour us and
those that do not alike, and it tests whether a number is stable and not what it means: an
admissible stratum whose margin is not positive, and an admissible stratum that is another one
restated, both establish nothing.
